%% kapa: exam=1 round=19 year=2008 period=2 form=A subject=economics subject_ko=경제원론 seq=39 number=39
%% source: Q-Net 감정평가사 기출문제 (cst00309) artlSeq=1001962 file=extracted/1008/1001962/2008년 제19회 감정평가사 1차시험 기출문제/2교시.hwp
%% conversion: kordoc 4.24.51 -> tools/gen1.py (auto); review=proofread-sonnet; flags=none
\begin{question}{39}
공유의 비극(tragedy of the commons)에 대한 설명 중 옳은 것은?
\begin{choices}
  \item 배제성과 경합성이 모두 있는 재화에서 발생한다.
  \item 배제성은 있으나 경합성이 없는 재화에서 발생한다.
  \item 배제성은 없으나 경합성이 있는 재화에서 발생한다.
  \item 배제성과 경합성이 모두 없는 재화에서 발생한다.
  \item 공공재와 자연독점의 경우에서 발생하는 현상이다.
\end{choices}
\end{question}
\answer{3}
